\section{Results}
\label{sec:results}

\begin{table}[t]
\centering
\footnotesize
\setlength{\tabcolsep}{3.5pt}
\resizebox{\linewidth}{!}{%
\begin{tabular}{lccc}
\toprule
Representation & Fold AUROC & Pooled AUROC & 95\% interval \\
\midrule
Handcrafted XGB & $0.644\pm0.086$ & 0.625 & [0.518, 0.744] \\
Frozen ESM3     & $0.761\pm0.131$ & \textbf{0.746} & [0.653, 0.837] \\
OAS LoRA ESM3   & $0.745\pm0.140$ & 0.726 & [0.626, 0.825] \\
\bottomrule
\end{tabular}
}
\caption{Classification performance for the selected handcrafted, frozen, and LoRA ESM3 configurations. A group bootstrap produces the intervals.}
\label{tab:main_classification}
\end{table}

\subsection{Classification}

Table~\ref{tab:main_classification} separates the configuration-selection score
(mean over folds) from the pooled OOF estimate and group-bootstrap interval.
Frozen ESM3 improves pooled AUROC from $0.625$ to $0.746$ and pooled AUPRC from
$0.742$ to $0.859$. On aligned predictions, the paired frozen-minus-handcrafted
difference is $+0.121$ AUROC ($[0.012,0.228]$) and $+0.117$ AUPRC
($[0.028,0.207]$). The intervals describe sensitivity to held-out lineage
composition; they do not remove model-selection optimism.

LoRA's selected classifier reaches pooled AUROC $0.726$, below frozen by
$-0.020$ with paired interval $[-0.096,0.059]$. The corresponding AUPRC
difference is $-0.031$ $[-0.075,0.024]$. These data do not establish a
downstream classification improvement from the single adaptation run. The
receiver operating characteristic and precision--recall curves for these
classifiers appear in Figure~\ref{fig:classification_curves}.

\begin{figure}[b]
    \centering
    \includegraphics[width=\linewidth]{../plots/lora/classification_roc_pr.png}
    \caption{Pooled OOF ROC and precision--recall curves for the selected classifiers.}
    \label{fig:classification_curves}
\end{figure}

The learned attention pool reaches mean-fold AUROC $0.702\pm0.100$, below the
best fixed-pooling frozen model. This is a negative capacity result on the
observed cohort, not evidence that learned pooling is generally inferior.

\subsection{Continuous breadth regression}
\label{sec:regression_results}

Regression uses the same 300 threshold-selected antibodies. As summarized in
Table~\ref{tab:regression}, frozen ESM3 reaches pooled Spearman $\rho=0.366$
($[0.126,0.585]$), Pearson $r=0.357$ ($[0.129,0.580]$), and $R^2=0.114$
($[-0.134,0.313]$). The length-only pooled Spearman floor is $0.004$. Thus the
selected representation contains ranking information, while its absolute breadth
predictions remain weak and sensitive to lineage composition. The scatter plots
of predicted versus true breadth values appear in
Figure~\ref{fig:regression_scatter}, illustrating the prediction variance.

Selecting LoRA configurations by mean-fold Spearman chooses random forest with
mean pooling, not the previously pooled-selected XGBoost configuration. Its
pooled Spearman is $0.319$ ($[0.059,0.561]$). The paired LoRA-minus-frozen
Spearman difference is $-0.047$ ($[-0.169,0.064]$).

\begin{table}[t]
\centering
\small
\setlength{\tabcolsep}{4pt}
\begin{tabular}{lccc}
\toprule
Representation & Fold $\rho$ & Pooled $\rho$ & 95\% interval \\
\midrule
Frozen ESM3   & $0.380\pm0.312$ & \textbf{0.366} & [0.126, 0.585] \\
OAS LoRA ESM3 & $0.357\pm0.320$ & 0.319 & [0.059, 0.561] \\
Length only   & --- & 0.004 & --- \\
\bottomrule
\end{tabular}
\caption{Continuous-breadth regression performance. Pooled OOF metrics and group-bootstrap intervals evaluate the selected configurations.}
\label{tab:regression}
\end{table}

\begin{figure}[b]
    \centering
    \includegraphics[width=\linewidth]{../plots/lora/regression_pred_vs_true.png}
    \caption{Pooled OOF predicted versus observed breadth for frozen and LoRA configurations.}
    \label{fig:regression_scatter}
\end{figure}

\subsection{Reconstruction dominates VAE interpolation}
\label{sec:vae_results}

The frozen-embedding VAE's best epoch has validation total loss $0.2599$,
reconstruction loss $0.1859$, and KL loss $14.80$. These values document the
saved run; they do not by themselves establish useful interpolation geometry.

In the inferred intermediate search (detailed in Table~\ref{tab:tier1}), raw
ESM3 interpolation is closer than the VAE path for six of seven VRC34 inferred
intermediates. Raw-line Mean Absolute Error (MAE) ranges from $8.51$ to $46.55$
(mean $20.77$), while VAE-path MAE ranges from $34.54$ to $44.30$ (mean
$37.47$). The exception is \texttt{pI3'}, for which the raw-line MAE is $46.55$
and the VAE MAE is $44.30$. The VAE-minus-\texttt{sl\_rec} difference averages
only $-0.0069$ (range $-0.0196$ to approximately zero), so the saved case study
does not demonstrate practically meaningful curvature beyond reconstruction.

The maturation midpoint test provides the cleanest direct midpoint comparison.
Across four ordered VRC34 triples, raw-line midpoint MAE ranges from $8.24$ to
$13.35$, versus $35.20$ to $37.95$ for the VAE midpoint. The per-residue VAE
bottleneck therefore adds much more error than the midpoint-control differences
under study.

\begin{table}[t]
\centering
\footnotesize
\setlength{\tabcolsep}{4pt}
\resizebox{\linewidth}{!}{%
\begin{tabular}{lrrr}
\toprule
Inferred intermediate & VAE MAE & Raw-line MAE & VAE--rec.\ line \\
\midrule
\texttt{pI1}      & 37.96 & 8.51  & $-0.0000$ \\
\texttt{pI2}      & 34.54 & 15.54 & $-0.0106$ \\
\texttt{pI2'}     & 38.28 & 32.48 & $-0.0196$ \\
\texttt{pI2 H33P} & 35.29 & 17.14 & $-0.0179$ \\
\texttt{pI3}      & 36.56 & 15.53 & $\phantom{-}0.0000$ \\
\texttt{pI3'}     & 44.30 & 46.55 & $-0.0000$ \\
\texttt{pI4}      & 35.38 & 9.63  & $-0.0000$ \\
\bottomrule
\end{tabular}
}
\caption{Exact embedding-space MAE per dimension/residue for VRC34 intermediates (inferred intermediate search). The reconstructed straight-line control provides the comparison for VAE performance.}
\label{tab:tier1}
\end{table}

\subsection{LoRA and descriptive maturation directionality}
\label{sec:lora_geometry}

For the maturation correlation analysis, mean Pearson correlation across the
four held-out lineages increases from $0.566$ for the frozen-embedding VAE to
$0.822$ for the LoRA-embedding VAE; mean Spearman increases from $0.384$ to
$0.692$. All four Pearson point estimates increase
(Figure~\ref{fig:directionality}), but sensitivity intervals are wide for small
lineages. For example, DH511 changes from $0.084$ to $0.409$, while its
approximate LoRA Pearson interval is $[-0.415,0.865]$.

This comparison is descriptive. It contrasts two separately trained, single-seed
VAEs, uses one LoRA adapter, covers four lineages, and lacks a raw-embedding
projection control. It therefore cannot establish that OAS adaptation caused
improved maturation organization.

The latent space geometry check does not provide supporting interpolation
evidence. Frozen VAE-minus-line MAE has mean $0.041$, standard deviation
$0.274$, and lineage-bootstrap interval $[-0.016,0.109]$; LoRA has mean $0.076$,
standard deviation $0.237$, and interval $[0.023,0.142]$. Positive differences
favor the reconstructed straight-line control.

\begin{figure}[b]
    \centering
    \includegraphics[width=\linewidth]{../plots/lora/vae_directionality_by_lineage.png}
    \caption{Pearson maturation correlation for four held-out lineages (maturation correlation analysis).}
    \label{fig:directionality}
\end{figure}
